\section{Central columns and descent of the dimension field}\label{sec:central-columns}

Let $\C$ be a self-dual rank-four complex fusion category,
$\mathcal S$ its sphericalization, and $\varepsilon$ its sign object.
Write $R=K_0(\C)$ and
$T=K_0(\mathcal S)/(\varepsilon+1)$, and put
$D=\operatorname{FPdim}(\C)$ and
$2E=\dim(\mathcal S)$. The ordinary and signed rational algebras
are commutative, as explained above. Let $L$ be the field generated
by the signed spherical dimensions. It is totally real and abelian.
Central induction gives $E/D\in L$, so $D\in L$.

We use the canonical normalized modular matrix
$S=s/(2E)$ of $Z(\mathcal S)$. For a one-dimensional character
$\chi$ of $K_0(\mathcal S)$, let $A_\chi$ be the corresponding
simple constituent of $I(\mathbf1)$. The central-induction theorem
and modular-character formula give
\begin{equation}\label{col:eq:central-column-character}
 [I(\mathbf1):A_\chi]=1,\qquad
 \dim A_\chi=\frac{2E}{f_\chi},\qquad
 \frac{S_{A_\chi,V}}{S_{A_\chi,0}}=\chi(F(V)).
\end{equation}
Here $F:Z(\mathcal S)\to\mathcal S$ is forgetful; see
\cite[Example~2.9 and Theorem~2.13]{Ostrik3}.

\begin{proposition}[Simultaneous Galois columns]\label{col:prop:simultaneous-columns}
Let $L$ be the signed dimension field and $\sigma\in\Gal(L/\Q)$.
Let $B_\sigma$ be the induction constituent associated with the signed
dimension character followed by $\sigma$. If $\chi$ is any ordinary
or signed character and $f_\chi$ its codegree in $R$ or $T$, then
\[
\chi(F(B_\sigma))=\frac{f_\chi}{\sigma(f_\chi)}.
\]
Consequently the codegrees determine both its folded and signed
forgetful classes, even if the rational algebras are products of fields.
Their coefficient vectors $w_\sigma,z_\sigma$ satisfy
\[
w_{\sigma,0}=z_{\sigma,0}=1,\qquad
w_{\sigma,i}\ge |z_{\sigma,i}|,\qquad
w_{\sigma,i}\equiv z_{\sigma,i}\pmod2.
\]
\end{proposition}
\begin{proof}
Let $A_\chi$ be the induction constituent for $\chi$. Its dimension
is $E/f_\chi$ and $S_{A_\chi,0}=1/(2f_\chi)$. The modular-character
formula gives
\[
S_{B_\sigma,A_\chi}
=\frac{\sigma(E/f_\chi)}{2\sigma(E)}
=\frac1{2\sigma(f_\chi)}.
\]
Symmetry of $S$ gives the asserted character value. All characters
together separate the commutative semisimple algebra, so these values
determine each class. If the two lift multiplicities at $X_i$ are
$m_i^+,m_i^-$, the two coefficients are their sum and difference.
Adjunction gives unit multiplicity one. The sign multiplicity is zero:
$I(\varepsilon)=\widehat\varepsilon\otimes I(1)$ has negative-dimensional
constituents, whereas the constituents of $I(1)$ have positive dimensions.
\end{proof}

\begin{proposition}[Positive Galois columns]\label{col:prop:positive-Galois-columns}
Suppose $[\Q(D):\Q]=4$. Identify $R$ with
its Frobenius--Perron image and $T$ with its spherical-dimension
image in the common field $K=L=\Q(D)$. For every
$\sigma\in\operatorname{Gal}(K/\Q)$, the element
\[
 u_\sigma=\frac{D}{\sigma(D)}
\]
belongs to $R$ and has nonnegative integral coefficients in its
distinguished basis. Write $w_\sigma$ for its coefficient vector and
$z_\sigma$ for the signed coefficient vector of
$v_\sigma=E/\sigma(E)\in T$. Then
\[
 w_{\sigma,0}=z_{\sigma,0}=1,\qquad
 w_{\sigma,i}\ge |z_{\sigma,i}|,\qquad
 w_{\sigma,i}\equiv z_{\sigma,i}\pmod2.
\]
\end{proposition}
\begin{proof}
Apply Proposition~\ref{col:prop:simultaneous-columns} to all ordinary
and signed characters. In a quartic field the dimension maps are
injective, so the folded and signed elements are respectively
$D/\sigma(D)$ and $E/\sigma(E)$. The coefficient assertions are
those of the proposition.
\end{proof}



\begin{proposition}[Descent of the dimension field]\label{col:prop:dimension-field-grading}
If $L\ne\Q(D)$, then $\C$ admits a nontrivial
faithful grading by a finite cyclic group.
\end{proposition}
\begin{proof}
Choose a nonidentity $\sigma\in\operatorname{Gal}(L/\Q(D))$.
Define $B_\sigma$ as in the preceding proof. The same
calculation applies without any injectivity assumption on the
dimension maps and gives
\[
 \operatorname{FPdim}(F(B_\sigma))=D/\sigma(D)=1.
\]
Adjunction gives $[F(B_\sigma):\mathbf1]=1$. Consequently
$F(B_\sigma)=\mathbf1$ in $\mathcal S$. The character defining
$B_\sigma$ differs from the dimension character, whose constituent
is the tensor unit. Thus $B_\sigma\ne\mathbf1$ in $Z(\mathcal S)$.

A central structure on the tensor unit is a monoidal natural
automorphism $\gamma$ of $\operatorname{Id}_{\mathcal S}$, and here
$\gamma$ is nonidentity. Its value at $\varepsilon$ is one.
Indeed the modular-character formula evaluated at the canonical
central lift $\widehat\varepsilon$ gives $-1$; on the other hand
this value is $\dim(\varepsilon)$ times the monodromy scalar,
namely $-\gamma_\varepsilon$. Hence
$\gamma_\varepsilon=1$.

The scalar values of $\gamma$ therefore agree on the two lifts of
every simple of $\C$. They descend to a nonidentity
monoidal natural automorphism of $\operatorname{Id}_{\C}$:
if an ordinary fusion coefficient is nonzero, one of its lifted
coefficients is nonzero and supplies the required product relation.
The finitely many scalar values form a subgroup of $\CC^\times$,
since tensor products give products and duals give inverses. This
finite group is cyclic and gives the asserted faithful grading.
\end{proof}

\begin{lemma}\label{col:lem:rank-four-grading}
A rank-four complex fusion category with a nontrivial faithful
grading is either weakly integral or has a nontrivial invertible
simple object.
\end{lemma}
\begin{proof}
The neutral component has rank one, two, or three. In rank one,
$X\otimes X^*$ is a sum of units for every simple $X$, and its
unit multiplicity is one. Thus every simple is invertible.

In rank two, the neutral component is pointed or has Fibonacci
fusion rules by \cite{OstrikRankTwo}. In the pointed case its global
dimension is two, and equality of the dimensions of homogeneous
components makes the whole category weakly integral. In the
Fibonacci case write $f^2=1+f$ and
$\phi=(1+\sqrt5)/2$. A component of rank one is impossible:
the integer multiplicity for multiplication by $f$ would satisfy
$m^2=1+m$. Thus there is just one other component, of rank two.
On its simples $X,Y$, multiplication by the self-dual object $f$
has a symmetric nonnegative integer matrix $M$ satisfying
$M^2=I+M$. Its eigenvalues are the two roots of $t^2-t-1$,
so its trace is one and determinant is minus one. Up to exchanging
$X,Y$,
\[
 M=\begin{pmatrix}0&1\\1&1\end{pmatrix},\qquad
 \operatorname{FPdim}(Y)=\phi\operatorname{FPdim}(X).
\]
Equality of component dimensions now gives
$(1+\phi^2)\operatorname{FPdim}(X)^2=1+\phi^2$.
Hence $X$ is a nontrivial invertible object.

In rank three, there is a single nonneutral simple $X$, and the
grading group has order two. Each neutral simple $Y$ satisfies
$Y\otimes X=m_YX$, so $\operatorname{FPdim}(Y)=m_Y$ is integral.
The two homogeneous components have equal Frobenius--Perron
dimension; hence the total dimension is an integer. See
\cite[Proposition~8.20]{ENO} for equality of component dimensions.
\end{proof}

\begin{corollary}\label{col:cor:quartic-order-casimir-degree}
If the rational Grothendieck algebra of a rank-four complex fusion
category is a quartic field, its Frobenius--Perron global dimension
has degree four.
\end{corollary}
\begin{proof}
Its Frobenius--Perron character identifies the rational algebra with
a real abelian quartic field $K$. Since this character is injective
and $\FPdim(X^*)=\FPdim(X)$, every simple object is self-dual.
Write $d_i$ for its simple
dimensions. If $D=\sum_i d_i^2$ is rational, all four formal codegrees
are $D$, so the reciprocal-codegree identity gives $D=4$. All four
$d_i$ are then one, contradicting injectivity of the character.

Otherwise, if $\Q(D)$ is a proper subfield of $K$, choose a
nonidentity $\sigma$ fixing $D$. Character spectral domination gives
$|\sigma(d_i)|\le d_i$. Equality of their square sums forces
$\sigma(d_i)=s_i d_i$, with $s_i\in\{\pm1\}$. Applying $\sigma$
to the dimension equations and comparing with the positive equations
shows that $s_i s_j=s_k$ whenever $N_{ij}^k>0$. Since the $d_i$
generate $K$, some $s_i$ is negative. Thus these signs define a
nontrivial grading. Lemma~\ref{col:lem:rank-four-grading} gives weak
integrality or a nontrivial invertible object. Weak integrality would
make $D$ rational. A nontrivial invertible is also impossible in a
totally real field algebra: it has finite order and thus equals
$1$ or $-1$, while its positive dimension is one. Both alternatives
are contradictory.
\end{proof}

\begin{corollary}\label{col:cor:rank-four-dimension-field-descent}
Every nonspherical self-dual rank-four complex fusion category
satisfies $L=\Q(D)$.
\end{corollary}
\begin{proof}
Otherwise the preceding proposition and lemma apply. Weakly integral
fusion categories are spherical by \cite[Proposition~8.24]{ENO}.
Theorem~\ref{inv:thm:rankfour-invertible-spherical} gives a spherical
structure in the remaining case, a contradiction.
\end{proof}

\subsection{The Haagerup ring}
\begin{theorem}\label{haag:thm:automatic-spherical}
Every complex fusion categorification of the rank-four Haagerup ring
admits a spherical structure.
\end{theorem}
\begin{proof}
In the tuple convention~\eqref{eq:tuple-convention}, this ring is
\[
 R_1=(0,0,1,1,1,1,1,1,2,2).
\]
Swapping $Y,Z$ gives the tuple $(0,1,0,1,1,1,2,2,1,1)$.
In the original order, the characters of $R_1$ are $(1,0,1,-1)$, $(1,1,-1,0)$, and
$(1,z-1,z,z+1)$ with $z^2=3z+1$. Consequently, writing
$\eta=(3+\sqrt{13})/2$, its FP codegree is
\[
 D=3\sqrt{13}\eta=\frac{39+9\sqrt{13}}2,
 \qquad DD'=117.
\]
Suppose the category has no spherical structure.
Corollary~\ref{col:cor:rank-four-dimension-field-descent} makes the
signed dimension field $K=\Q(\sqrt{13})$.
Select its larger dimension codegree $E\ge E'$.
Then $E<D$, while central induction makes $EE'/117$ a positive integer.

The fundamental unit of $K$ is $\eta$, of norm $-1$.
Indeed a positive unit of norm $-1$ has the form
$(a+b\sqrt{13})/2$, with $a,b>0$ and $a^2+4=13b^2$;
therefore $a\ge3$, with equality at $\eta$.
A smaller positive unit of norm one would produce a smaller
norm-minus-one unit by division into $\eta$.
Because $E$ is a $d$-number, $E/E'$ is a totally positive unit,
so $E/E'=\eta^{2h}$ for an integer $h\ge0$. Set $y=E\eta^{-h}$.
Since $\eta'=-\eta^{-1}$, one has $y'=(-1)^h y$.
Thus $y\in\Q$ if $h$ is even and
$y\in\Q\sqrt{13}$ if $h$ is odd.
The element $y$ is integral, because $\eta$ is an integral unit.
As $\mathcal O_K=\ZZ[(1+\sqrt{13})/2]$, the two intersections
are respectively $\ZZ$ and $\ZZ\sqrt{13}$.
Hence the divisibility $117\mid EE'$ gives precisely
\[
 E=3s\sqrt{13}\eta^h\quad(h\ge1\text{ odd}),
 \qquad\text{or}\qquad
 E=39s\eta^h\quad(h\ge0\text{ even}),
 \qquad s\in\ZZ_{>0}.
\]
For the first form, its norm is $13$ times the square of its rational
coefficient, forcing that coefficient to be divisible by three.
For the second, the rational coefficient has square divisible by
$117=9\cdot13$, forcing divisibility by $39$.
The first form satisfies $E\ge D$; the second satisfies
$E\ge39>D$. Both contradict properness.
\end{proof}

